% 第4章：本章小结。

\section{本章小结}
\label{sec:ls-operations-summary}

本章按定位一个数所需的索引个数，依次处理向量、矩阵与张量，并对每种运算交代输入条件、输出形状和可用的规则。向量的加法与数乘合成线性组合；内积把两个向量配成一个数，范数用四条公理统一各种误差口径，由内积给出的欧氏范数还满足平行四边形恒等式。

矩阵乘积来自两次加权求和的复合，满足结合律与分配律，却不满足交换律与消去律；逆矩阵撤销一次乘法，矩阵幂记录反复作用。迹、行列式与矩阵范数把矩阵读成一个数：行列式是有向体积的缩放倍数，诱导范数是最坏情形的放大倍数。张量把索引推广到任意多个：重排只改变元素位置，张量积保留全部槽位，缩约沿指定的轴求和，内积、矩阵乘积与迹都是缩约的特例。

本章只把这些数组当作具体的数表。第~\ref{chap:linear-representation}章接着追问它们记录的是什么：同一个对象换一组方向后坐标怎样改变，矩阵怎样记录线性映射，一般空间中的夹角与正交怎样讨论。矩阵反复作用的走向与误差的放大留给第~\ref{chap:linear-analysis}章，方程怎样求解留给第~\ref{chap:linear-equations-and-optimization}章。
